\chapter{Evaluación VLM/ICL}
\label{chap:vlm}

Este capítulo define la línea VLM/ICL del estudio: modelos de visión y lenguaje con aprendizaje en contexto. Las tablas no incorporan cifras para no anticipar resultados VLM; el capítulo fija el diseño experimental, las condiciones, los controles, el formato de salida y la lectura comparativa.

\newcommand{\vlmpending}{--}

\section{Objetivo de la línea VLM/ICL}

La línea CNN ha establecido la base supervisada: con \(K \leq 32\), ResNet18 aprende en el simulador, pero no alcanza un rendimiento fiable sobre los datos cedidos por el HUCA. El protocolo VLM/ICL evalúa la hipótesis complementaria: si los ejemplos etiquetados son demasiado escasos para actualizar pesos de forma fiable, quizá puedan ser más útiles como demostraciones dentro del contexto de un modelo multimodal preentrenado.

La comparación se formula con la misma entrada visual, los mismos pacientes, los mismos presupuestos \(K\), la misma separación leave-one-out y las mismas métricas por paciente. La diferencia experimental es una sola: la CNN usa los ejemplos para entrenar; el VLM los usa para condicionar la inferencia.

\section{Modelos evaluados}

El protocolo contempla dos modelos abiertos:

\begin{itemize}
  \item \textbf{Gemma 4}. Modelo multimodal generalista con entrada de texto e imagen \cite{gemma4card}. Actúa como referencia abierta no especializada. Su papel es medir cuánto puede aportar un VLM general con contexto visual y textual amplio.
  \item \textbf{MedGemma 1.5}. Modelo multimodal médico basado en la familia Gemma y adaptado con modalidades clínicas \cite{medgemma15}. Actúa como referencia médica. Su papel es medir si el preentrenamiento biomédico general aporta ventaja en ECG, aun cuando la documentación del modelo no identifique el ECG como modalidad principal.
\end{itemize}

Ambos modelos se sirven localmente mediante vLLM siempre que el hardware lo permita. Esta decisión evita depender de APIs externas y permite registrar versión de pesos, plantilla de conversación, resolución visual, temperatura y parámetros de generación.

\section{Presupuestos y demostraciones}

Los presupuestos son \(K=\{0,2,4,8,16,32\}\). El caso \(K=0\) mide comportamiento \emph{zero-shot}: el modelo recibe instrucciones y la imagen objetivo, pero ninguna demostración etiquetada. Los valores positivos de \(K\) miden aprendizaje en contexto con ejemplos imagen-respuesta.

Para que la comparación con la CNN sea directa, las demostraciones se seleccionan a nivel de paciente y nunca incluyen el paciente reservado para test. Cuando un paciente aporta varias derivaciones, la demostración conserva la misma lógica de agregación que el resto del estudio: V1, V2 y V3 son evidencia visual del mismo paciente, no muestras independientes para filtrar entre entrenamiento y prueba.

El esquema general de esta construcción se resume en la figura~\ref{fig:icl_context_window}. En este capítulo se concreta cómo se seleccionan las demostraciones, qué campos debe producir el modelo y cómo se agregan las respuestas por paciente.

\section{Formato de salida}

El formato depende de la tarea VLM. En la tarea morfológica, el JSON estricto contiene los mismos hallazgos que utiliza la CNN. Para cada imagen V1, V2 o V3, la respuesta mínima es:

\begin{lstlisting}
{
  "RBBB": true,
  "ST_ELEVATION": true,
  "T_WAVE_INVERSION": true
}
\end{lstlisting}

La etiqueta clínica de paciente se deriva después mediante la misma regla AND usada en el resto del TFG: solo se marca sospecha Brugada si los tres hallazgos están presentes tras agregar las derivaciones. Esta tarea permite estudiar transferencia desde demostraciones morfológicas sintéticas hacia la cohorte del HUCA.

La segunda tarea es clínica directa. En ella el VLM recibe una única derivación fija del paciente, por defecto V2, y responde:

\begin{lstlisting}
{
  "clinical_brugada": true
}
\end{lstlisting}

La tarea clínica directa usa \(K\) pacientes reales de la cohorte cedida por el HUCA como demostraciones en contexto y predice otro paciente de esa misma cohorte. Para \(K \geq 2\), la selección de contexto fuerza ejemplos normales y Brugada siempre que el pliegue lo permita. La derivación se fija de antemano, por defecto V2, de modo que ambos modelos ven exactamente las mismas imágenes. Con ello se mide si el VLM aprovecha pocos ejemplos clínicos reales sin actualizar pesos, una pregunta complementaria a la transferencia morfológica desde simulador.

El análisis principal no interpreta texto libre. Una respuesta que no pueda analizarse como JSON, con claves ausentes, tipos incorrectos o valores fuera del esquema, se registra como fallo de formato. Esta tasa de fallos se reporta junto con las métricas clínicas porque forma parte del comportamiento real de un VLM.

\section{Condiciones experimentales}

La tabla \ref{tab:vlm_conditions} fija las condiciones de la campaña. Las tres primeras condiciones miden rendimiento; las dos últimas son controles de causalidad multimodal.

\begin{table}[H]
  \centering
  \caption{Condiciones experimentales VLM/ICL.}
  \label{tab:vlm_conditions}
  \small
  \begin{tabularx}{\textwidth}{p{3.2cm}p{2.2cm}X}
    \toprule
    Condición & \(K\) & Propósito \\
    \midrule
    Zero-shot & 0 & Medir conocimiento previo sin demostraciones. \\
    ICL normal & 2, 4, 8, 16, 32 & Medir aprendizaje en contexto con pares imagen-etiqueta correctos. \\
    ICL balanceado & 2, 4, 8, 16, 32 & Mantener proporción positiva/negativa controlada cuando el pliegue lo permita. \\
    Etiquetas permutadas & 8, 16, 32 & Romper la correspondencia imagen-etiqueta para detectar uso espurio del contexto. \\
    Sin imagen de apoyo & 8, 16, 32 & Mantener etiquetas de demostración, pero retirar las imágenes de apoyo. \\
    \bottomrule
  \end{tabularx}
\end{table}

\section{Agregación por paciente}

En la tarea morfológica, cada predicción se obtiene por imagen y se agrega a nivel de paciente para conservar la misma unidad de evaluación que en la CNN. Si el modelo produce una respuesta válida para V1, V2 y V3, se calcula una decisión de paciente por voto mayoritario o por media de probabilidades si se activa la confianza numérica. En la tarea clínica directa, la consulta contiene una sola imagen del paciente; la salida \texttt{clinical\_brugada} se evalúa directamente contra la referencia clínica de la cohorte. Si alguna respuesta produce JSON inválido, se conserva el fallo en la auditoría. El análisis principal contabiliza los fallos como errores; un análisis secundario puede informar métricas sobre respuestas válidas, pero nunca sustituye al análisis principal.

\section{Trazabilidad de la evaluación}

La evaluación VLM/ICL conserva la misma separación por paciente que la línea CNN. Para cada condición se registran el modelo, la tarea, el presupuesto \(K\), la semilla, las imágenes de consulta, las demostraciones incluidas, la respuesta estructurada y la validez del formato. Esta información permite reconstruir métricas por paciente, matrices de confusión y comparaciones por presupuesto sin mezclar dominios ni pacientes.

\section{Tablas de resultados VLM/ICL}

Las tablas siguientes mantienen las celdas numéricas sin valor para no adelantar resultados VLM. En esta sección, \(\vlmpending\) identifica una celda sin cifra reportada.

\begin{table}[H]
  \centering
  \caption{Tabla de resultados VLM/ICL para Gemma 4 sobre los datos cedidos por el HUCA.}
  \label{tab:vlm_gemma4_results}
  \small
  \begin{tabular}{rrrrrrr}
    \toprule
    \(K\) & BA & F1 & Sens. & Esp. & JSON inválido & Latencia media \\
    \midrule
    0  & \vlmpending & \vlmpending & \vlmpending & \vlmpending & \vlmpending & \vlmpending \\
    2  & \vlmpending & \vlmpending & \vlmpending & \vlmpending & \vlmpending & \vlmpending \\
    4  & \vlmpending & \vlmpending & \vlmpending & \vlmpending & \vlmpending & \vlmpending \\
    8  & \vlmpending & \vlmpending & \vlmpending & \vlmpending & \vlmpending & \vlmpending \\
    16 & \vlmpending & \vlmpending & \vlmpending & \vlmpending & \vlmpending & \vlmpending \\
    32 & \vlmpending & \vlmpending & \vlmpending & \vlmpending & \vlmpending & \vlmpending \\
    \bottomrule
  \end{tabular}
\end{table}

\begin{table}[H]
  \centering
  \caption{Tabla de resultados VLM/ICL para MedGemma 1.5 sobre los datos cedidos por el HUCA.}
  \label{tab:vlm_medgemma_results}
  \small
  \begin{tabular}{rrrrrrr}
    \toprule
    \(K\) & BA & F1 & Sens. & Esp. & JSON inválido & Latencia media \\
    \midrule
    0  & \vlmpending & \vlmpending & \vlmpending & \vlmpending & \vlmpending & \vlmpending \\
    2  & \vlmpending & \vlmpending & \vlmpending & \vlmpending & \vlmpending & \vlmpending \\
    4  & \vlmpending & \vlmpending & \vlmpending & \vlmpending & \vlmpending & \vlmpending \\
    8  & \vlmpending & \vlmpending & \vlmpending & \vlmpending & \vlmpending & \vlmpending \\
    16 & \vlmpending & \vlmpending & \vlmpending & \vlmpending & \vlmpending & \vlmpending \\
    32 & \vlmpending & \vlmpending & \vlmpending & \vlmpending & \vlmpending & \vlmpending \\
    \bottomrule
  \end{tabular}
\end{table}

\begin{table}[H]
  \centering
  \caption{Controles VLM/ICL para verificar uso multimodal real.}
  \label{tab:vlm_controls_results}
  \small
  \begin{tabularx}{\textwidth}{p{3.1cm}p{2.2cm}XXXX}
    \toprule
    Modelo & Condición & \(K\) & BA & Caída frente a normal & Lectura \\
    \midrule
    Gemma 4 & Etiquetas permutadas & 16 & \vlmpending & \vlmpending & \vlmpending \\
    Gemma 4 & Sin imagen de apoyo & 16 & \vlmpending & \vlmpending & \vlmpending \\
    MedGemma 1.5 & Etiquetas permutadas & 16 & \vlmpending & \vlmpending & \vlmpending \\
    MedGemma 1.5 & Sin imagen de apoyo & 16 & \vlmpending & \vlmpending & \vlmpending \\
    Gemma 4 & Etiquetas permutadas & 32 & \vlmpending & \vlmpending & \vlmpending \\
    Gemma 4 & Sin imagen de apoyo & 32 & \vlmpending & \vlmpending & \vlmpending \\
    MedGemma 1.5 & Etiquetas permutadas & 32 & \vlmpending & \vlmpending & \vlmpending \\
    MedGemma 1.5 & Sin imagen de apoyo & 32 & \vlmpending & \vlmpending & \vlmpending \\
    \bottomrule
  \end{tabularx}
\end{table}

\section{Matriz de comparación}

La comparación se lee como una curva de decisión. La CNN aporta los valores de referencia y la columna VLM conserva el mismo espacio de lectura, sin introducir cifras no verificadas.

\begin{table}[H]
  \centering
  \caption{Matriz de decisión CNN frente a VLM/ICL.}
  \label{tab:cnn_vlm_decision_matrix}
  \footnotesize
  \begin{tabularx}{\textwidth}{@{}lcccX@{}}
    \toprule
    \(K\) & CNN base & CNN comp. & VLM & Lectura \\
    \midrule
    2  & 0,509 & -- & \vlmpending & Zona favorable a ICL. \\
    4  & 0,509 & -- & \vlmpending & Zona favorable a ICL. \\
    8  & 0,490 & -- & \vlmpending & Zona de transición. \\
    16 & 0,516 & 0,524 (CORAL) & \vlmpending & Comparación abierta. \\
    32 & 0,499 & 0,549 (MMD) & \vlmpending & Comparación abierta. \\
    \bottomrule
  \end{tabularx}
\end{table}

\section{Figuras de comparación VLM/ICL}

\begin{figure}[H]
  \centering
  \fcolorbox{MidGray}{PaperGray}{
    \parbox[c][4.2cm][c]{0.88\textwidth}{
      \centering
      \textbf{Comparación CNN vs VLM/ICL}\\[0.7em]
      \small Exactitud equilibrada por \(K\), con CNN base, mejor comprobación complementaria, Gemma 4 y MedGemma 1.5.
    }
  }
  \caption{Espacio de comparación de exactitud equilibrada entre CNN y VLM/ICL.}
  \label{fig:vlm_cnn_comparison_reserved}
\end{figure}

\begin{figure}[H]
  \centering
  \fcolorbox{MidGray}{PaperGray}{
    \parbox[c][4.2cm][c]{0.88\textwidth}{
      \centering
      \textbf{Matrices de confusión VLM}\\[0.7em]
      \small Matrices por modelo y \(K\), incluyendo conteo de respuestas JSON inválidas.
    }
  }
  \caption{Espacio para matrices de confusión de Gemma 4 y MedGemma 1.5.}
  \label{fig:vlm_confusion_reserved}
\end{figure}

\section{Criterio de lectura}

La hipótesis principal es que ICL será más competitivo en \(K\) bajo, donde entrenar pesos con pocos pacientes resulta frágil. Si el VLM supera a la CNN en esa zona, el resultado apoyará el uso de demostraciones en contexto para tareas raras con pocas etiquetas. Si la CNN alcanza o supera al VLM cuando \(K\) aumenta, aparecerá un punto de cruce práctico: por debajo conviene ICL; por encima conviene entrenamiento supervisado. Si ninguno de los dos enfoques supera un rendimiento aceptable, el resultado también será informativo, porque indicará que la tarea requiere más datos, mejores anotaciones morfológicas o modelos específicamente preentrenados en ECG.

En todos los casos, la lectura clínica se mantiene constante. La salida del modelo solo prioriza revisión experta dentro del circuito de Brugada; no confirma ni descarta el síndrome.
